% Part: incompleteness
% Chapter: incompleteness-provability
% Section: provability-conditions

\documentclass[../../../include/open-logic-section]{subfiles}

\begin{document}

\olfileid[pt]{inc}{inp}{prc}

\olsection{As condições de \usetoken{S}{derivability} para $\Th{PA}$}

A aritmética de Peano, ou $\Th{PA}$, é a teoria que estende $\Th{Q}$ com
axiomas de indução para todas as !!{formula}s. Em outras palavras,
acrescentamos a $\Th{Q}$ axiomas da forma
\[
(!A(0) \land \lforall[x][(!A(x) \lif !A(x'))]) \lif \lforall[x][!A(x)]
\]
para toda !!{formula}~$!A$. Observe que isso é, na verdade, um
\emph{esquema}, isto é, uma coleção infinita de axiomas (e verifica-se que
$\Th{PA}$ {\em não} é finitamente axiomatizável). Mas, como podemos
determinar efetivamente se uma cadeia de símbolos é ou não uma instância de
um axioma de indução, o conjunto de axiomas de $\Th{PA}$ é computável.
$\Th{PA}$ é uma teoria muito mais robusta que~$\Th{Q}$. Por exemplo, usando
indução da maneira usual, podemos provar facilmente que a adição e a
multiplicação são comutativas. De fato, a maior parte dos argumentos
finitários da teoria dos números e da combinatória pode ser desenvolvida
em~$\Th{PA}$.

Como $\Th{PA}$ é axiomatizada computavelmente, o predicado de
!!{derivability} $\Prf[\Th{PA}](x,y)$ é computável e, portanto,
representado em~$\Th{Q}$ (e, assim, em~$\Th{PA}$). Como antes, tomaremos
$\OPrf[\Th{PA}](x,y)$ para denotar a fórmula que representa a relação. Seja
$\OProv[\Th{PA}](y)$ a fórmula
$\lexists[x][\Prf[\Th{PA}](x,y)]$, que intuitivamente diz: ``$y$ é
!!{derivable} a partir dos axiomas de $\Th{PA}$.'' Precisamos de um pouco
mais que os axiomas de $\Th{Q}$ porque precisamos saber que a teoria usada
é forte o bastante para !!{derive} alguns fatos básicos sobre esse
predicado de !!{derivability}. De fato, precisamos dos fatos seguintes:
\begin{enumerate}
\item[P1.] Se $\Th{PA} \Proves !A$, então $\Th{PA} \Proves
  \OProv[\Th{PA}](\gn{!A})$.
\item[P2.] Para todas as !!{formula}s $!A$ e $!B$,
  \[
  \Th{PA} \Proves \OProv[\Th{PA}](\gn{!A \lif !B}) \lif
  (\OProv[\Th{PA}](\gn{!A}) \lif \OProv[\Th{PA}](\gn{!B})).
  \]
\item[P3.] Para toda !!{formula}~$!A$,
  \[
  \Th{PA} \Proves \OProv[\Th{PA}](\gn{!A})
  \lif \OProv[\Th{PA}](\gn{\OProv[\Th{PA}](\gn{!A})}).
  \]
\end{enumerate}
A única maneira de verificar que essas três propriedades valem é descrever
cuidadosamente a !!{formula} $\OProv[\Th{PA}](y)$ e usar os axiomas de
$\Th{PA}$ para descrever as !!{derivation}s formais pertinentes. As
condições (1) e~(2) são fáceis; é realmente a condição~(3) que exige
trabalho. (Pense no tipo de trabalho que ela envolve \dots) Desenvolver os
detalhes seria tedioso e pouco interessante, por isso pediremos aqui que
você aceite pela confiança que $\Th{PA}$ tem as três propriedades listadas
acima. Uma escolha razoável de $\OProv[\Th{PA}](y)$ também satisfará
\begin{enumerate}
\item[P4.] Se $\Th{PA} \Proves \OProv[\Th{PA}](\gn{!A})$, então
  $\Th{PA} \Proves !A$.
\end{enumerate}
Mas não precisaremos desse fato.

\begin{digress}
Aliás, G\"odel foi preguiçoso da mesma maneira que estamos sendo agora. No
fim do artigo de 1931, ele esboça a prova do segundo teorema da incompletude
e promete os detalhes em um artigo posterior. Nunca chegou a escrevê-los;
como todos os que compreendiam o argumento acreditavam que ele podia ser
desenvolvido, ele não precisou preencher os detalhes.
\end{digress}

\end{document}
